% Part: second-order-logic
% Chapter: syntax-and-semantics
% Section: language-of-sol

\documentclass[../../../include/open-logic-section]{subfiles}

\begin{document}

\olfileid{sol}{syn}{lan}

\olsection{دوجے درجے دی منطق دی زبان}

پہلے درجے دی منطق وانگ دوجے درجے دی منطق دیاں عبارتاں اک بنیادی لفظ نامے توں بندیاں نیں، جیہدے وچ \emph{!!{variable}s}، \emph{!!{constant}s}، \emph{!!{predicate}s} تے کدے \emph{!!{function}s} ہوندے نیں۔ ایہناں توں، منطقی جوڑاں، مسوّراں تے قوساں تے کامیاں ورگیاں وقفے دیاں نشانیاں نال رلا کے، \emph{ٹرم} تے \emph{!!{formula}s} بندے نیں۔ فرق ایہ اے کہ شیواں دے متغیراں توں علاوہ دوجے درجے دی منطق وچ تعلقاں تے فنکشناں دے متغیر وی نیں، تے ایہ اوہناں اُتے مسوّر لاؤن دی کھُل دیندی اے۔ سو دوجے درجے دی منطق دیاں منطقی نشانیاں پہلے درجے دی منطق والیاں نیں، تے نال ایہ:

\begin{enumerate}
\item اک !!{denumerable}s سیٹ دوجے درجے دے تعلقاں دے !!{variable}s دا، ہر تھاں-گنتی~$n$ لئی: $\Obj V_0^n$، $\Obj V_1^n$، $\Obj V_2^n$، \dots
\item اک !!{denumerable}s سیٹ دوجے درجے دے فنکشناں دے !!{variable}s دا:
  $\Obj u_0^n$، $\Obj u_1^n$، $\Obj u_2^n$، \dots
\end{enumerate}

پہلے درجے دی منطق وچ !!{identity}~$\eq$ عام طور اُتے شامل ہوندی اے۔ پہلے درجے دی منطق وچ زبان~$\Lang{L}$ دیاں غیر منطقی نشانیاں کوئی دلچسپ گل ظاہر کرن لئی ضروری نیں۔ بے شک !!{sentences} وی نیں جیہڑے کوئی غیر منطقی نشانی نہیں ورتدے، پر صرف~$\eq$ نال کوئی دلچسپ گل کہنا اوکھا اے۔ دوجے درجے دی منطق وچ تعلقاں تے فنکشناں دے متغیر بے حد نیں، ایس لئی جیہڑی وی گل پہلے درجے دی زبان وچ کہہ سکدے آں، اوہ ایتھے غیر منطقی نشانیاں دی خاص رسد توں بناں وی کہہ سکدے آں۔

\end{document}
